\documentclass[../../main.tex]{subfiles}
\begin{document}

\begin{center}
  {\fontsize{48pt}{58pt}\selectfont \textbf{1994}\normalsize\textbf{年度}}
\end{center}

\begin{table}[htb]
  \centering
  \caption{出題テーマと難易度}
  \begin{tabular}{|c|p{9cm}|c|}
    \hline
    \textbf{問題番号} & \textbf{テーマ} & \textbf{難易度} \\
    \hline
    第1問 & 行列の可換条件と平方根 & \\
    \hline
    第2問 & フィボナッチ型数列の剰余と和 & \\
    \hline
    第3問 & 正四面体の切断平面と円の交点 & \\
    \hline
    第4問 & 角の二等分線で分ける線分比の極値 & \\
    \hline
    第5問 & 札の交換と確率漸化式 & \\
    \hline
    第6問 & 媒介変数表示の曲線の長さと極限 & \\
    \hline
  \end{tabular}
\end{table}

\begin{table}[htb]
  \centering
  \caption{解答}
  \renewcommand{\arraystretch}{1.4}
  \begin{tabular}{|c|c|p{9cm}|}
    \hline
    \textbf{問題番号} & \textbf{小問} & \textbf{答え} \\
    \hline
    第1問 & (1) &  \\
    \cline{2-3}
    & (2) &  \\
    \cline{2-3}
    & (3) &  \\
    \hline
    第2問 & (1) & $b_0,\ldots,b_9=1,2,0,2,2,1,0,1,1,2$ \\
    \cline{2-3}
    & (2) & $c_{n+8}=c_n+9$（証明問題） \\
    \cline{2-3}
    & (3) & $n+1\le c_n\le\dfrac{3}{2}(n+1)$（証明問題） \\
    \hline
    第3問 & -- & $0<s,t<1$，$t\ge\dfrac{4s}{1+3s}$ \\
    \hline
    第4問 & -- & 最大値$\dfrac{3\sqrt{5}}{5}$，最小値$\dfrac{\sqrt{5}}{3}$ \\
    \hline
    第5問 & (1) & $a_n=\dfrac{b_{n-1}+2c_{n-1}}{3}$，$b_n=\dfrac{a_{n-1}+2b_{n-1}}{3}$，$c_n=\dfrac{2a_{n-1}+c_{n-1}}{3}$ \\
    \cline{2-3}
    & (2) & $a_{2k}=\dfrac{1+2(1/3)^k}{3}$，$a_{2k+1}=\dfrac{1-(1/3)^k}{3}$ \\
    \hline
    第6問 & (1) & $L=16$ \\
    \cline{2-3}
    & (2) & $\lim_{n\to\infty}\sqrt{n}\,\theta_n=4$ \\
    \hline
  \end{tabular}
\end{table}

\paragraph{試験形態}
6問．試験時間は要確認．

\end{document}
